% JETSPIN evaporation model documentation
% Yarin, Koombhongse and Reneker, J. Appl. Phys. 89, 3018--3026 (2001)

\section{Solvent evaporation and solidification\label{sec:Evaporation}}

JETSPIN implements the solvent-evaporation model introduced by Yarin,
Koombhongse and Reneker \cite{yarin2001bending}.  The dimensional mass-transfer
correlation is
\begin{equation}
 Sh=0.495\,Re^{1/3}Sc^{1/2},
 \qquad
 Re=\frac{2a_i|\mathbf{v}_i|}{\nu_a},
 \qquad
 Sc=\frac{\nu_a}{D_a},
 \label{eq:evap-sherwood}
\end{equation}
where $a_i$ is the local jet radius, $\nu_a$ is the kinematic viscosity
of air, and $D_a$ is the diffusion coefficient of solvent vapour in air.
For a discrete jet element of length $l_i$ and volume
$V_i=\pi a_i^2l_i$, the corresponding evaporation equation is
\begin{equation}
 \frac{dV_i}{dt}=
 -0.495\,D_a Re_i^{1/3}Sc^{1/2}
 \left[c_{s,\mathrm{eq}}(T)-c_{s,\infty}\right]\pi l_i .
 \label{eq:evap-volume-dimensional}
\end{equation}
This is the discrete form of Eq.~(34) of Yarin et al.
\cite{yarin2001bending}.

The ambient value used by JETSPIN is obtained from the relative humidity
$RH$ according to
\begin{equation}
 c_{s,\infty}=RH\,c_{s,\mathrm{eq}}(T).
 \label{eq:evap-humidity}
\end{equation}
If the saturation vapour concentration is not supplied explicitly,
JETSPIN evaluates it from the temperature using the correlation of
Seaver, Galloway and Manuccia employed by Yarin et al.

\subsection{JETSPIN nondimensionalization}

The evaporation equations are integrated together with the equations of
motion in the standard JETSPIN dimensionless variables.  With
\begin{equation}
 t_0=\frac{\mu_0}{G_0},
 \qquad
 L_0=\sqrt{\frac{q_0^2}{\pi a_0^2G_0}},
\end{equation}
JETSPIN uses
\begin{equation}
 \bar t=\frac{t}{t_0},\quad
 \bar l_i=\frac{l_i}{L_0},\quad
 \bar{\mathbf v}_i=\frac{t_0}{L_0}\mathbf v_i,\quad
 \bar V_i=\frac{V_i}{L_0^3},
\end{equation}
and
\begin{equation}
 \bar\nu_a=\nu_a\frac{t_0}{L_0^2},
 \qquad
 \bar D_a=D_a\frac{t_0}{L_0^2}.
\end{equation}
Both $Re$ and $Sc$ are invariant under this rescaling.  Consequently,
Eq.~\ref{eq:evap-volume-dimensional} becomes
\begin{equation}
 \frac{d\bar V_i}{d\bar t}=
 -0.495\,\bar D_a Re_i^{1/3}Sc^{1/2}
 \left[c_{s,\mathrm{eq}}(T)-c_{s,\infty}\right]\pi\bar l_i,
 \label{eq:evap-volume-dimensionless}
\end{equation}
which is the expression evaluated in \texttt{eom\_ev\_mod.f90}.
The arrays \texttt{jetvl} and \texttt{jetve} contain, respectively,
the reference volume $\bar V_i^0$ of the material element before solvent
loss and its instantaneous volume $\bar V_i$ after evaporation.

\subsection{Polymer concentration and rheological solidification}

Conservation of polymer mass gives
\begin{equation}
 c_{p,i}=c_{p,0}\frac{V_i^0}{V_i}.
 \label{eq:evap-cp}
\end{equation}
The viscosity is updated according to Eqs.~(37) and (41) of Yarin et al.:
\begin{equation}
 \frac{\mu_i}{\mu_0}=
 10^{B\left(c_{p,i}^{m}-c_{p,0}^{m}\right)}.
 \label{eq:evap-viscosity}
\end{equation}
For the Yarin--Koombhongse--Reneker model the relaxation time obeys
\begin{equation}
 \frac{\theta_i}{\theta_0}=\frac{c_{p,i}}{c_{p,0}},
 \label{eq:evap-relaxation}
\end{equation}
and, since $G_i=\mu_i/\theta_i$,
\begin{equation}
 \frac{G_i}{G_0}=
 \frac{10^{B\left(c_{p,i}^{m}-c_{p,0}^{m}\right)}}
      {c_{p,i}/c_{p,0}}.
 \label{eq:evap-elastic}
\end{equation}
JETSPIN retains a generalized exponent $t_{\mathrm{ev}}$ through
$\theta_i/\theta_0=(c_{p,i}/c_{p,0})^{t_{\mathrm{ev}}}$; setting
\texttt{evaporation tconstant 1} gives exactly Eq.~\ref{eq:evap-relaxation}
and is the Yarin-2001 choice.  The default value is one.

For Maxwell rheology, with $\bar\sigma_i=\sigma_i/G_0$, the resulting
constitutive equation is
\begin{equation}
 \frac{d\bar\sigma_i}{d\bar t}=
 \frac{1}{\theta_i/\theta_0}
 \left[
 \frac{\mu_i}{\mu_0}\frac{1}{\bar l_i}
 \frac{d\bar l_i}{d\bar t}-\bar\sigma_i
 \right],
 \label{eq:evap-maxwell}
\end{equation}
which is Eq.~(39) of Yarin et al. written in the JETSPIN scaling.

\subsection{Evaporation cutoff}

Yarin et al. stop the evaporation calculation when the solvent mass
fraction
\begin{equation}
 c_{s,i}=1-c_{p,i}
\end{equation}
reaches $c_{s,i}=0.1$.  Using Eq.~\ref{eq:evap-cp}, the equivalent
condition on the volume represented by a bead is
\begin{equation}
 \frac{V_i}{V_i^0}=\frac{c_{p,0}}{1-0.1}=\frac{c_{p,0}}{0.9}.
 \label{eq:evap-cutoff}
\end{equation}
The JETSPIN integrators enforce Eq.~\ref{eq:evap-cutoff} at every
intermediate and final integration stage.  Once this limit is reached,
$V_i$, $c_{p,i}$, $\mu_i$, $\theta_i$, and $G_i$ remain constant, as in
the solidification cutoff adopted in Ref.~\cite{yarin2001bending}.

\subsection{Input directives}

The following directives activate and parameterize evaporation:
\begin{description}
 \item[\texttt{evaporation yes}] activate solvent evaporation and the
 concentration-dependent rheology.
 \item[\texttt{evaporation polymer frac} $f$] initial polymer mass
 fraction $c_{p,0}$.
 \item[\texttt{evaporation airviscosity} $f$] kinematic viscosity of
 air $\nu_a$ in $\mathrm{cm^2/s}$.
 \item[\texttt{evaporation temperature} $f$] temperature in kelvin.
 \item[\texttt{evaporation umidity} $f$] relative humidity $RH$ in the
 interval $[0,1]$.  The historical spelling \texttt{umidity} is retained
 for input-file compatibility.
 \item[\texttt{evaporation diffusivity} $f$] vapour diffusion coefficient
 $D_a$ in $\mathrm{cm^2/s}$.
 \item[\texttt{evaporation bconstant} $f$] parameter $B$ in
 Eq.~\ref{eq:evap-viscosity}.
 \item[\texttt{evaporation mconstant} $f$] exponent $m$ in
 Eq.~\ref{eq:evap-viscosity}.
 \item[\texttt{evaporation tconstant} $f$] exponent
 $t_{\mathrm{ev}}$ in the generalized relaxation-time law; use $f=1$
 to reproduce Yarin et al. (2001).
\end{description}

The example in \texttt{examples/input-7/input.dat} uses Maxwell rheology
and \texttt{evaporation tconstant 1}, and therefore follows the
Yarin-2001 rheological law.
